\documentclass[11pt,a4paper]{article}

% ---------- Packages ----------
\usepackage[utf8]{inputenc}
\usepackage[T1]{fontenc}
\usepackage{lmodern}
\usepackage[margin=2.5cm]{geometry}
\usepackage{amsmath,amssymb,amsthm}
\usepackage{graphicx}
\usepackage{booktabs}
\usepackage{caption}
\usepackage{subcaption}
\usepackage[numbers,sort&compress]{natbib}
\usepackage{xcolor}
\usepackage[colorlinks=true,linkcolor=blue,citecolor=blue,urlcolor=blue]{hyperref}
\usepackage[capitalise]{cleveref}

% ---------- Theorem environments ----------
\newtheorem{theorem}{Theorem}
\newtheorem{lemma}[theorem]{Lemma}
\theoremstyle{definition}
\newtheorem{definition}{Definition}

% ---------- Handy macros ----------
\newcommand{\todo}[1]{\textcolor{red}{[TODO: #1]}}

% ---------- Title ----------
\title{Paper Title Goes Here}
\author{
  First Author\thanks{Corresponding author: \texttt{email@example.com}} \\
  Affiliation \\
  \and
  Second Author \\
  Affiliation
}
\date{\today}

\begin{document}

\maketitle

\begin{abstract}
\todo{One paragraph (150--250 words): context, problem, method, key results, implications.}
\end{abstract}

\noindent\textbf{Keywords:} keyword one, keyword two, keyword three

% ======================================================
\section{Introduction}
\label{sec:intro}
\todo{Motivation and background.}

\todo{Problem statement and research gap.}

\todo{Contributions of this paper:}
\begin{itemize}
  \item Contribution 1.
  \item Contribution 2.
  \item Contribution 3.
\end{itemize}

The rest of the paper is organised as follows: \cref{sec:related} reviews related work,
\cref{sec:method} describes the method, \cref{sec:experiments} reports experiments,
\cref{sec:discussion} discusses the findings, and \cref{sec:conclusion} concludes.

% ======================================================
\section{Related Work}
\label{sec:related}
\todo{Group prior work by theme and contrast it with this paper, e.g.\ \citet{example2024}.}

% ======================================================
\section{Method}
\label{sec:method}

\subsection{Problem Formulation}
\begin{definition}
\todo{Formal definition of the problem.}
\end{definition}

An example equation:
\begin{equation}
  \mathcal{L}(\theta) = \frac{1}{N}\sum_{i=1}^{N} \ell\big(f_\theta(x_i), y_i\big).
  \label{eq:loss}
\end{equation}

\subsection{Proposed Approach}
\todo{Describe the approach; refer to \cref{eq:loss} and \cref{fig:overview}.}

\begin{figure}[t]
  \centering
  \fbox{\parbox[c][4cm][c]{0.8\linewidth}{\centering Placeholder figure}}
  % \includegraphics[width=0.8\linewidth]{figures/overview.pdf}
  \caption{Overview of the proposed approach.}
  \label{fig:overview}
\end{figure}

% ======================================================
\section{Experiments}
\label{sec:experiments}

\subsection{Setup}
\todo{Datasets, baselines, metrics, implementation details.}

\subsection{Results}
\todo{Main results; see \cref{tab:results}.}

\begin{table}[t]
  \centering
  \caption{Main results. Best values in bold.}
  \label{tab:results}
  \begin{tabular}{lccc}
    \toprule
    Method        & Metric A       & Metric B       & Metric C \\
    \midrule
    Baseline 1    & 0.00           & 0.00           & 0.00 \\
    Baseline 2    & 0.00           & 0.00           & 0.00 \\
    Ours          & \textbf{0.00}  & \textbf{0.00}  & \textbf{0.00} \\
    \bottomrule
  \end{tabular}
\end{table}

\subsection{Ablation Study}
\todo{Effect of each component.}

% ======================================================
\section{Discussion}
\label{sec:discussion}
\todo{Interpretation of results, limitations, threats to validity.}

% ======================================================
\section{Conclusion}
\label{sec:conclusion}
\todo{Summary of contributions and future work.}

% ======================================================
\section*{Acknowledgements}
\todo{Funding and acknowledgements.}

\bibliographystyle{plainnat}
\bibliography{references}

\appendix
\section{Additional Results}
\label{app:extra}
\todo{Supplementary material.}

\end{document}
